\chapter{Antipsychotic Medicines}

\section*{Definition and Overview}
Antipsychotic medicines are drugs used primarily to treat \textbf{psychosis} -- the group of symptoms that includes hallucinations, delusions, and disorganised thinking. They are the cornerstone of treatment for schizophrenia and for mania in bipolar disorder. They are also used, at lower doses and in certain circumstances, for severe agitation, tics, some obsessive-compulsive symptoms, and as an add-on treatment for depression that has not responded to other medicines.

There are two broad groups. \textbf{First-generation} (typical) antipsychotics, such as haloperidol and chlorpromazine, were introduced in the 1950s. \textbf{Second-generation} (atypical) antipsychotics, such as risperidone, olanzapine, quetiapine, and aripiprazole, are now more commonly used. They are available as tablets, liquids, and long-acting depot injections given every few weeks. Antipsychotics do not cure psychotic illness, but they control symptoms effectively and dramatically reduce the risk of relapse when taken continuously.

\section*{Epidemiology}
Schizophrenia affects around 1 in 100 people worldwide and typically begins in late adolescence or early adulthood, slightly earlier in men than in women. Bipolar disorder is similarly common. Together, these conditions are among the leading causes of disability in young adults, and antipsychotics are the mainstay of their treatment.

Antipsychotic prescribing has widened beyond these core diagnoses in recent decades. They are increasingly prescribed, often off-label, for severe anxiety, resistant depression, and agitation associated with dementia in older people. Because long-term treatment is common and these drugs cause weight gain and metabolic changes, people taking antipsychotics carry a higher burden of diabetes, cardiovascular disease, and premature mortality, making metabolic monitoring an essential part of care.

\section*{Aetiology and Pathophysiology}
Antipsychotics act mainly by blocking \textbf{dopamine D2 receptors} in the brain. Psychosis is linked to excess dopamine activity in the mesolimbic pathways that process thinking and perception, so blocking these receptors dampens hallucinations and delusions. The same mechanism, however, produces the characteristic side effects of the class:

\begin{itemize}
    \item \textbf{Dopamine blockade:} reduces psychotic symptoms, but also causes movement side effects (extrapyramidal symptoms) and a rise in the hormone prolactin
    \item \textbf{Serotonin receptor blockade:} a feature of most atypical drugs that lessens movement side effects but increases weight gain and metabolic disturbance
    \item \textbf{Histamine, muscarinic, and alpha-adrenergic effects:} responsible for drowsiness, dry mouth, constipation, and postural low blood pressure
    \item \textbf{Long-term receptor adaptation:} sustained treatment is far more effective at preventing relapse than episodic use, in part because the brain stabilises on a steady drug level
\end{itemize}

\section*{Clinical Features}
\begin{itemize}
    \item \textbf{Intended effect:} reduced hallucinations and delusions, clearer thinking, and less agitation or mania, typically over days to weeks
    \item \textbf{Sedation and drowsiness:} common, especially with olanzapine and quetiapine, and most noticeable in the first weeks
    \item \textbf{Weight gain and metabolic changes:} increased appetite, weight gain, and rises in blood sugar and cholesterol
    \item \textbf{Movement effects:} restlessness (akathisia), muscle stiffness, tremor, drooling, and slow stiff movements, particularly with first-generation drugs
    \item \textbf{Hormonal effects:} breast swelling, breast milk production (galactorrhoea), and changes to periods or sexual function from raised prolactin
    \item \textbf{Tardive dyskinesia:} with long-term use, involuntary movements of the face, tongue, and limbs that may persist
    \item \textbf{Autonomic effects:} dry mouth, blurred vision, constipation, and light-headedness on standing
\end{itemize}

\section*{Differential Diagnosis}
When symptoms are not improving or new problems appear, the possibilities to weigh up include:

\begin{itemize}
    \item \textbf{Inadequate treatment:} too low a dose, the wrong medicine, or missed doses versus genuinely treatment-resistant illness
    \item \textbf{Side effects of the medicine} mistaken for worsening of the underlying condition
    \item \textbf{Substance-induced psychosis:} from alcohol, cannabis, stimulants, or withdrawal
    \item \textbf{Medical causes of psychosis or confusion:} delirium from infection, electrolyte disturbance, thyroid disease, or neurological conditions
    \item \textbf{Neuroleptic malignant syndrome:} fever, severe muscle rigidity, and confusion, which must be distinguished from infection
    \item \textbf{Relapse of the illness} versus a new medical problem
\end{itemize}

\section*{When to Seek Medical Care}
See a doctor promptly for distressing movement symptoms, rapid weight gain, or signs of diabetes such as excessive thirst, frequent urination, and weight loss. A medication review is also needed before stopping or changing treatment.

\textbf{Call 000 (or local emergency services) for:}
\begin{itemize}
    \item High fever with severe muscle stiffness, confusion, sweating, and rapid heart rate -- signs of neuroleptic malignant syndrome
    \item Thoughts of suicide, self-harm, or plans to end your life
    \item Sudden weakness, facial droop, or difficulty speaking, especially in an older person taking these medicines
    \item Fainting, collapse, severe agitation, or an inability to care for oneself
    \item Swelling of the face or throat or difficulty breathing (allergic reaction)
\end{itemize}

\section*{Diagnosis and Investigations}
\begin{itemize}
    \item \textbf{Psychiatric assessment:} confirming the diagnosis, the indication for the medicine, and the suitability of long-term treatment
    \item \textbf{Baseline physical checks:} weight, waist circumference, blood pressure, blood glucose, and lipid profile before starting
    \item \textbf{Regular monitoring:} repeated weight and metabolic testing (at least every six months) during treatment
    \item \textbf{ECG:} baseline and periodic heart rhythm checks for drugs associated with QT prolongation
    \item \textbf{Prolactin measurement:} if symptoms suggest raised prolactin
    \item \textbf{Adherence review:} for depot injections, a check that doses are received and drug levels are adequate
    \item \textbf{Movement disorder screening:} regular examination using standardised scales, including for tardive dyskinesia
\end{itemize}

\section*{Management}
\begin{itemize}
    \item \textbf{Choice of drug:} individualised according to the condition, past response, and the side-effect profile a person can tolerate
    \item \textbf{Start low and go slow:} doses are built up gradually, reviewed regularly, and kept at the lowest effective level
    \item \textbf{Long-acting depot injections:} an excellent option when daily tablets are difficult to maintain, with dosing every two to four weeks
    \item \textbf{Psychological and social support:} cognitive behavioural therapy for psychosis, family support, and vocational rehabilitation substantially improve outcomes
    \item \textbf{Managing side effects:} dietary and activity measures for weight, and specific medicines or dose adjustment for movement effects
    \item \textbf{Smoking:} smoking induces enzymes that break down some antipsychotics, so dose adjustments may be needed when smoking habits change
    \item \textbf{Shared care:} treatment usually initiated and supervised by a psychiatrist, with regular review by the general practitioner and mental health team
\end{itemize}

\section*{Complications}
\begin{itemize}
    \item Weight gain and metabolic syndrome, raising the risk of type 2 diabetes and cardiovascular disease
    \item Extrapyramidal movement disorders, including tardive dyskinesia, which may be persistent and disabling
    \item Neuroleptic malignant syndrome, a rare but life-threatening emergency
    \item Raised prolactin, with effects on breasts, periods, and sexual function, and an increased risk of bone thinning over time
    \item Sedation, postural hypotension, and falls, especially in older people
    \item An increased risk of stroke and death when antipsychotics are used for behavioural symptoms in older people with dementia
    \item QTc prolongation and rare cardiac arrhythmias with some agents
    \item Relapse of the illness if medicines are stopped suddenly
\end{itemize}

\section*{Prognosis and Outcomes}
Antipsychotics do not cure schizophrenia or bipolar disorder, but they are very effective at controlling symptoms and preventing relapse in most people, and they make it possible to study, work, and live an active life. Long-term, regular treatment is considerably more effective than intermittent or crisis-only use.

The course of psychotic illness varies widely: some people have a single episode, while others experience repeated relapses, most often when medication is stopped. Weight gain and metabolic side effects need active, ongoing management because they strongly influence long-term physical health. With good adherence, family support, and rehabilitation, a substantial proportion of people achieve sustained recovery.

\section*{Prevention and Self-Care}
\begin{itemize}
    \item Take medicines reliably and attend depot appointments; stopping suddenly is the single most common cause of relapse
    \item Follow a healthy diet and get regular exercise to counter weight gain
    \item Attend scheduled weight, blood sugar, blood pressure, and heart checks
    \item Avoid alcohol and recreational drugs, which can worsen psychosis and interfere with treatment
    \item Learn your early warning signs of relapse and agree on an action plan with your team
    \item Stay connected with family, support workers, and community services
    \item Report side effects early rather than quietly stopping the medicine
    \item Use sleep routines and stress management, since poor sleep is a common trigger for relapse
\end{itemize}

\section*{Sources and Further Reading}
The following organisations provide reliable, up-to-date information on antipsychotic medicines and psychotic illness for patients, students, and health professionals.

\begin{itemize}
    \item NHS. \textit{Information on antipsychotics and treatment of psychosis and schizophrenia.} https://www.nhs.uk/conditions/
    \item NICE. \textit{Guidelines on schizophrenia and psychosis, including medication.} https://www.nice.org.uk/
    \item Mayo Clinic. \textit{Overview of antipsychotics, uses, and side effects.} https://www.mayoclinic.org/
    \item MedlinePlus. \textit{Health information on antipsychotic medicines and their effects.} https://medlineplus.gov/
    \item World Health Organization. \textit{Resources on mental health and psychotic disorders.} https://www.who.int/
    \item MSD Manual. \textit{Professional overview of schizophrenia and antipsychotic drugs.} https://www.msdmanuals.com/
\end{itemize}
